\documentclass[12pt]{book}
\usepackage{precalc1}
\setcounter{chapter}{3}
\begin{document}

\chapter{Rational Functions}

A rational function is a ratio of two polynomials,
\[
f(x)=\frac{p(x)}{q(x)},
\]
and its graph is governed by the interplay between the two. Where the numerator
$p$ is zero the function is zero; where the denominator $q$ is zero the function is
undefined and the graph runs off toward a vertical line. The relative sizes of the
two degrees decide what happens at the far ends. As with polynomials, the factored
form of numerator and denominator displays the information that fixes the shape, and
the leading terms settle the ends.

\section{Zeros and Vertical Asymptotes}

The numerator and denominator each contribute their own roots, and the two kinds
behave oppositely. A root of the numerator (that is not also a root of the
denominator) makes the whole fraction zero, so the graph crosses or touches the axis
there exactly as a polynomial would. A root of the denominator makes the fraction
undefined; near such a point the denominator is tiny while the numerator is not, so
the quotient grows without bound. The graph approaches a vertical line called a
vertical asymptote.

\begin{definition}{Zero and vertical asymptote}
For $f(x)=p(x)/q(x)$ in lowest terms, a \emph{zero} is an input where $p(x)=0$,
producing $f(x)=0$. A \emph{vertical asymptote} is a vertical line $x=a$ where
$q(a)=0$; as $x\to a$ the values of $f$ increase or decrease without bound.
\end{definition}

\begin{keyidea}
Numerator roots pull the graph \emph{to} the axis (zeros). Denominator roots push
the graph \emph{away} along a vertical line (asymptotes). The factored form of
$p/q$ lists both at once.
\end{keyidea}

\begin{visual}
At a numerator root the fraction reads (small)$/$(ordinary), which is small: the
graph sits near the $x$-axis. At a denominator root the fraction reads
(ordinary)$/$(small), which is enormous: the graph races up or down beside the
vertical line $x=a$. The two roots are mirror situations, one driving the value to
zero, the other to infinity.
\end{visual}

\begin{example}[locating zeros and asymptotes]
For
\[
f(x)=\frac{(x-2)(x+1)}{(x-3)(x+1)}
\]
the numerator is zero at $x=2$ and $x=-1$; the denominator is zero at $x=3$ and
$x=-1$. The factor $(x+1)$ is common to both, so it cancels: it is neither a true
zero nor an asymptote, but a single missing point (a hole) at $x=-1$. After
cancellation the function is $f(x)=\dfrac{x-2}{x-3}$ with a genuine zero at $x=2$ and
a vertical asymptote at $x=3$.
\end{example}

\begin{pitfall}
A factor common to numerator and denominator does not create an asymptote. It
cancels, leaving a hole, a single point removed from the graph. Always reduce to
lowest terms before declaring zeros and asymptotes; the leftover denominator roots
are the asymptotes, the leftover numerator roots are the zeros.
\end{pitfall}

\subsection{Sign analysis across the features}

Zeros and vertical asymptotes are the only places a rational function in lowest terms
can change sign, for the same reason as with polynomials: the value can switch sign
only by passing through zero (a numerator root) or through infinity (a denominator
root). Marking both kinds on a number line and testing one point per interval gives
the sign on each interval, just as before.

\begin{example}[sign chart for a rational function]
Take $f(x)=\dfrac{x-2}{x-3}$. The feature points are the zero $x=2$ and the
asymptote $x=3$.
\begin{center}
\renewcommand{\arraystretch}{1.3}
\begin{tabular}{@{}lcccc@{}}
\toprule
Interval & test $x$ & $x-2$ & $x-3$ & sign of $f$\\
\midrule
$(-\infty,2)$ & $0$ & $-$ & $-$ & $+$\\
$(2,3)$       & $2.5$ & $+$ & $-$ & $-$\\
$(3,\infty)$  & $4$ & $+$ & $+$ & $+$\\
\bottomrule
\end{tabular}
\end{center}
The graph is above the axis left of $2$, dips below between $2$ and $3$, and returns
above the axis past the asymptote at $3$.
\end{example}

\section{End Behavior: Comparing Degrees}

Far from the origin the leading terms of numerator and denominator dominate all
others, so the end behavior of $p(x)/q(x)$ matches that of the ratio of leading
terms. Writing $p(x)=a_m x^m+\cdots$ and $q(x)=b_n x^n+\cdots$, the comparison of the
degrees $m$ and $n$ decides the outcome.

\begin{keyidea}
\textbf{End behavior of} $f(x)=\dfrac{a_m x^m+\cdots}{b_n x^n+\cdots}$
\begin{itemize}\setlength{\itemsep}{2pt}
  \item \emph{$m<n$ (proper, numerator smaller).} The ends approach the horizontal
        asymptote $y=0$.
  \item \emph{$m=n$ (equal degrees).} The ends approach the horizontal asymptote
        $y=\dfrac{a_m}{b_n}$, the ratio of leading coefficients.
  \item \emph{$m>n$ (improper, numerator larger).} There is no horizontal asymptote.
        Polynomial long division produces a polynomial part the ends follow; when
        $m=n+1$ that part is a line, a \emph{slant asymptote}.
\end{itemize}
\end{keyidea}

\begin{example}[horizontal asymptote from equal degrees]
For $f(x)=\dfrac{3x^2-1}{2x^2+x+5}$ the degrees are equal at $2$. The ends approach
$y=\tfrac{3}{2}$, the ratio of the leading coefficients $3$ and $2$. The lower-degree
terms affect the middle of the graph but vanish in the comparison at the ends.
\end{example}

\begin{example}[horizontal asymptote at zero]
For $f(x)=\dfrac{4x+1}{x^2-9}$ the numerator degree $1$ is below the denominator
degree $2$, so the ends approach $y=0$. The denominator factors as $(x-3)(x+3)$,
giving vertical asymptotes at $x=3$ and $x=-3$; the single numerator root
$x=-\tfrac14$ is the only zero.
\end{example}

\section{Improper Rational Functions and Long Division}

When the numerator degree is at least the denominator degree, the function is
\emph{improper}, and the honest description of its ends comes from polynomial long
division. Dividing $p$ by $q$ rewrites the function as a polynomial quotient plus a
proper remainder term:
\[
\frac{p(x)}{q(x)}=Q(x)+\frac{R(x)}{q(x)},\qquad \deg R<\deg q.
\]
Far from the origin the remainder term $R(x)/q(x)$ is proper, so it decays to zero,
and the graph approaches the polynomial $Q(x)$. The quotient $Q$ is the end-behavior
curve.

\begin{keyidea}
For an improper rational function, long division gives
$f(x)=Q(x)+\dfrac{R(x)}{q(x)}$ with the remainder term proper. The ends of the graph
hug the quotient $Q(x)$, because the proper part dies off. When $\deg p=\deg q+1$ the
quotient is a line, and that line is the slant asymptote.
\end{keyidea}

\begin{example}[a slant asymptote by long division]
Analyze $f(x)=\dfrac{x^2+x-2}{x-3}$.

The numerator degree $2$ exceeds the denominator degree $1$ by one, so expect a
slant asymptote. Divide $x^2+x-2$ by $x-3$:
\[
\begin{array}{r}
x+4\phantom{xxxxx}\\
x-3\,\big)\,\overline{x^2+x-2\phantom{x}}\\
\underline{x^2-3x\phantom{xxxx}}\\
4x-2\phantom{x}\\
\underline{4x-12}\\
10\phantom{x}
\end{array}
\]
So
\[
f(x)=x+4+\frac{10}{x-3}.
\]
The remainder term $\dfrac{10}{x-3}$ is proper and decays to zero at the ends, so the
graph approaches the line $y=x+4$. That line is the slant asymptote. The factored
numerator $x^2+x-2=(x+2)(x-1)$ gives zeros at $x=-2$ and $x=1$, and the denominator
gives a vertical asymptote at $x=3$.
\end{example}

\begin{center}
\begin{tikzpicture}[xscale=0.42,yscale=0.30]
  \draw[->] (-5.6,0)--(12.8,0) node[right] {$x$};
  \draw[->] (0,-8.6)--(0,18.6) node[above] {$y$};
  % vertical asymptote x=3
  \draw[pcpurple,dashed,thick] (3,-8.4) -- (3,18.4);
  \node[pcpurple] at (4.2,17.6) {\scriptsize $x=3$};
  % slant asymptote y=x+4
  \draw[pcgreen,dashed,thick,domain=-5.4:12.6] plot (\x,{\x+4});
  \node[pcgreen] at (10.8,16.4) {\scriptsize $y=x+4$};
  % left branch  (x<3)
  \draw[pcblue,thick,domain=-5.4:2.32,smooth,samples=160]
    plot (\x,{\x+4+10/(\x-3)});
  % right branch (x>3)
  \draw[pcblue,thick,domain=4.0:12.6,smooth,samples=160]
    plot (\x,{\x+4+10/(\x-3)});
  % zeros
  \fill[pcred] (-2,0) circle (3.2pt) node[below left] {\scriptsize $-2$};
  \fill[pcred] (1,0) circle (3.2pt) node[below right] {\scriptsize $1$};
\end{tikzpicture}
\end{center}

\begin{pitfall}
A slant asymptote appears only when the numerator degree is exactly one more than the
denominator degree. If it exceeds by two or more, the quotient $Q(x)$ is a parabola
or higher curve, and the ends follow that curve, not a line. There is no horizontal
asymptote in any improper case.
\end{pitfall}

\begin{example}[improper with equal-after-division check]
For $f(x)=\dfrac{2x^3}{x^2+1}$ the numerator degree exceeds the denominator by one.
Long division gives
\[
\frac{2x^3}{x^2+1}=2x-\frac{2x}{x^2+1},
\]
so the slant asymptote is $y=2x$ and the proper remainder $-\dfrac{2x}{x^2+1}$ decays
at the ends. The only zero is $x=0$, and the denominator $x^2+1$ is never zero, so
there are no vertical asymptotes.
\end{example}

\section{Assembling the Graph}

A faithful sketch of a rational function in lowest terms combines four readings, each
taken from the factored form or the long division.

\begin{keyidea}
\textbf{Sketching a rational function.}
\begin{enumerate}\setlength{\itemsep}{2pt}
  \item \emph{Reduce} to lowest terms; note any holes from cancelled factors.
  \item \emph{Zeros:} numerator roots, with crossing or touching set by multiplicity.
  \item \emph{Vertical asymptotes:} denominator roots; the graph escapes to $\pm\infty$
        beside each.
  \item \emph{End behavior:} compare degrees for a horizontal asymptote, or divide for
        a slant or polynomial asymptote.
  \item Use a sign chart over zeros and asymptotes to keep the curve on the correct
        side of the axis between features.
\end{enumerate}
\end{keyidea}

\begin{example}[a complete analysis]
Analyze $f(x)=\dfrac{x^2-1}{x^2-4}$.

\emph{Reduce.} Numerator $(x-1)(x+1)$, denominator $(x-2)(x+2)$; no common factors,
so no holes.

\emph{Zeros.} $x=1$ and $x=-1$, both simple, so the graph crosses there.

\emph{Vertical asymptotes.} $x=2$ and $x=-2$.

\emph{End behavior.} Equal degrees, leading coefficients both $1$, so the horizontal
asymptote is $y=1$.

\emph{Sign chart.} Feature points $-2,-1,1,2$ split the line into five intervals;
testing one point in each places the curve above the axis on the outer intervals
(approaching $y=1$), with the expected escapes beside $x=\pm2$.
\end{example}

\begin{center}
\begin{tikzpicture}[xscale=0.62,yscale=0.34]
  \draw[->] (-5.4,0)--(5.4,0) node[right] {$x$};
  \draw[->] (0,-6.4)--(0,6.4) node[above] {$y$};
  % vertical asymptotes
  \draw[pcpurple,dashed,thick] (-2,-6.2)--(-2,6.2);
  \draw[pcpurple,dashed,thick] (2,-6.2)--(2,6.2);
  \node[pcpurple] at (-2.0,6.9) {\scriptsize $x=-2$};
  \node[pcpurple] at (2.0,6.9) {\scriptsize $x=2$};
  % horizontal asymptote y=1
  \draw[pcgreen,dashed,thick] (-5.3,1)--(5.3,1);
  \node[pcgreen] at (4.7,1.7) {\scriptsize $y=1$};
  % outer-left branch  x in [-5,-2.45]  (above HA, rising to +inf)
  \draw[pcblue,thick,domain=-5:-2.45,smooth,samples=120]
    plot (\x,{(\x*\x-1)/(\x*\x-4)});
  % inner-left branch  x in [-1.86,1.86] continuous through the two zeros and the bowl,
  % but dips to -inf near +-2, so split: [-1.86,-1]...[-1,1]...[1,1.86] is one smooth curve
  \draw[pcblue,thick,domain=-1.86:1.86,smooth,samples=200]
    plot (\x,{(\x*\x-1)/(\x*\x-4)});
  % outer-right branch x in [2.45,5]
  \draw[pcblue,thick,domain=2.45:5,smooth,samples=120]
    plot (\x,{(\x*\x-1)/(\x*\x-4)});
  % zeros
  \fill[pcred] (-1,0) circle (3pt) node[above left] {\scriptsize $-1$};
  \fill[pcred] (1,0) circle (3pt) node[above right] {\scriptsize $1$};
\end{tikzpicture}
\end{center}

\demo{Build a rational function by placing numerator zeros and vertical asymptotes,
each with a multiplicity, in factored form, then fit a hidden target. The fit score
weighs zeros, asymptotes, multiplicities, and overall agreement, and the current
degree ratio reports the end behavior. Compare how a numerator root drives the curve
to the axis while a denominator root drives it to a vertical
line.}{https://bobovski66.github.io/Precalculus/demos/rational\_functions.html}

\section{Exercises}

\begin{exercises}
\item For $f(x)=\dfrac{(x-4)(x+2)}{(x-1)(x+5)}$, state the zeros and the vertical
asymptotes, and confirm there are no common factors.

\item Identify the hole, the zeros, and the vertical asymptotes of
$f(x)=\dfrac{(x-3)(x+2)}{(x+2)(x-5)}$ after reducing to lowest terms.

\item Explain, in terms of (small)$/$(ordinary) versus (ordinary)$/$(small), why a
numerator root produces a zero while a denominator root produces a vertical
asymptote.

\item Give the horizontal asymptote, if any, of each:
(a) $\dfrac{2x+1}{x-7}$; (b) $\dfrac{5x^2-3}{x^2+4}$; (c) $\dfrac{x+6}{x^3-1}$;
(d) $\dfrac{x^3}{x+2}$.

\item For $f(x)=\dfrac{6x^2+x}{3x^2-12}$, find the horizontal asymptote from the
leading coefficients, the zeros, and the vertical asymptotes.

\item Use long division to write $f(x)=\dfrac{x^2-5x+6}{x-1}$ in the form
$Q(x)+\dfrac{R}{x-1}$, and state the slant asymptote.

\item Find the slant asymptote of $f(x)=\dfrac{2x^2+3x-1}{x+2}$ by long division, and
give the zeros and the vertical asymptote.

\item Explain why an improper rational function never has a horizontal asymptote, and
state the condition under which its end-behavior curve is a straight line.

\item Build a complete sign chart for $f(x)=\dfrac{x-1}{(x+2)(x-3)}$ over its zeros
and vertical asymptotes, recording the sign of each factor in each interval.

\item Carry out a full analysis of $f(x)=\dfrac{x^2-4}{x^2-1}$ (reduce, zeros,
vertical asymptotes, end behavior, sign chart) and sketch a graph consistent with all
of it.

\item A rational function has zeros at $x=0$ and $x=3$, vertical asymptotes at
$x=-1$ and $x=2$, and a horizontal asymptote at $y=1$. Construct a factored formula
meeting all four conditions, and explain how you chose the degrees.

\item Open the Zeros and Asymptotes Lab. Build a rational function with at least two
vertical asymptotes and one repeated numerator zero, record its factored form, the
degree ratio, and the fit score against a target, and write one sentence on how the
end behavior changed when you altered the degree of the numerator.
\end{exercises}

\end{document}
